% UPI-TRM ICLR 2026 Workshop Poster
% Portrait, 61cm x 91cm (24in x 36in)
\documentclass[20pt, a1paper, portrait]{tikzposter}

\setlength{\paperwidth}{24in}
\setlength{\paperheight}{36in}
\makeatletter
\setlength{\TP@visibletextwidth}{22in}
\setlength{\TP@visibletextheight}{34.5in}
\setlength{\TP@blockverticalspace}{8mm}
\makeatother

% Disable tikzposter watermark
\tikzposterlatexaffectionproofoff

% ── Packages ──
\usepackage{amsmath, amssymb, mathtools}
\usepackage{booktabs}
\usepackage{enumitem}
\usepackage[T1]{fontenc}
\usepackage{lmodern}
\usepackage{microtype}
\usepackage{xcolor}
\usepackage{array}
\usepackage{multirow}
\usepackage{tikz}
\usetikzlibrary{positioning, arrows.meta, calc, decorations.pathreplacing}
\usepackage{pgfplots}
\pgfplotsset{compat=1.18}

% Suppress minor line-breaking issues
\emergencystretch=2em

% ── Macros ──
\newcommand{\R}{\mathbb{R}}
\newcommand{\E}{\mathbb{E}}
\newcommand{\norm}[1]{\left\lVert #1 \right\rVert}

% ── Colors: navy/blue/orange palette (color-blind safe) ──
\definecolor{TRMBlue}{HTML}{3498DB}
\definecolor{KnobGray}{HTML}{7F8C8D}
\definecolor{AccentGreen}{HTML}{27AE60}
\definecolor{AccentOrange}{HTML}{E67E22}
\definecolor{HeaderBG}{HTML}{1A2740}
\definecolor{EqBoxBG}{HTML}{EBF5FB}
\definecolor{ResultGold}{HTML}{F39C12}
\definecolor{NavyDark}{HTML}{2C3E50}
\definecolor{HeroBG}{HTML}{FDF2E9}
\definecolor{HeroBorder}{HTML}{E67E22}
% Color-blind safe bar chart colors
\definecolor{BarNoContr}{HTML}{E74C3C}
\definecolor{BarContr}{HTML}{2980B9}

\colorlet{EqBoxBGcolor}{EqBoxBG}
\colorlet{TRMBlueBorder}{TRMBlue}
\colorlet{OrangeBorder}{AccentOrange}
\colorlet{GreenLight}{AccentGreen!8!white}
\colorlet{GreenBorder}{AccentGreen}
\colorlet{BlueLight}{TRMBlue!8!white}
\colorlet{OrangeLight}{AccentOrange!8!white}
\colorlet{HeaderLight}{HeaderBG!8!white}
\colorlet{HeaderBorder}{HeaderBG}
\colorlet{HeroBGcolor}{HeroBG}
\colorlet{HeroBordercolor}{HeroBorder}

% ── Theme ──
\usetheme{Default}

\definecolorstyle{UPIStyle}{
  \definecolor{colorOne}{HTML}{1A2740}
  \definecolor{colorTwo}{HTML}{2C3E50}
  \definecolor{colorThree}{HTML}{3498DB}
}{
  \colorlet{backgroundcolor}{white}
  \colorlet{framecolor}{colorTwo}
  \colorlet{titlefgcolor}{white}
  \colorlet{titlebgcolor}{colorOne}
  \colorlet{blocktitlebgcolor}{colorTwo}
  \colorlet{blocktitlefgcolor}{white}
  \colorlet{blockbodybgcolor}{white}
  \colorlet{blockbodyfgcolor}{black}
}
\usecolorstyle{UPIStyle}

\defineblockstyle{UPIBlock}{
  titlewidthscale=1.0, bodywidthscale=1.0, titleleft,
  titleoffsetx=0pt, titleoffsety=0pt,
  bodyoffsetx=0pt, bodyoffsety=0pt, bodyverticalshift=0pt,
  roundedcorners=8, linewidth=1.5pt,
  titleinnersep=4mm, bodyinnersep=5mm
}{
  \draw[color=framecolor, fill=blocktitlebgcolor,
        rounded corners=\blockroundedcorners]
    (blocktitle.south west) rectangle (blocktitle.north east);
  \draw[color=framecolor, fill=blockbodybgcolor,
        rounded corners=\blockroundedcorners]
    (blockbody.south west) rectangle (blockbody.north east);
}
\useblockstyle{UPIBlock}

% ── Title ──
\title{\textsf{\textbf{Unrolled Policy Iteration for Tiny Recursive Models}}}
\author{\textsf{Bahram Behzadian$^{*1}$\enspace
  Brett Daley$^{*1}$\enspace
  Gopeshh Subbaraj$^{2}$\enspace
  Houssam Nassif$^{1}$}}
\institute{\textsf{$^{1}$Meta\qquad
  $^{2}$Mila -- Quebec AI Institute\qquad
  $^{*}$Equal contribution}}

\begin{document}

\maketitle[width=21in, innersep=8mm, titletoblockverticalspace=16mm]

% ═══════════════════════════════════
% HERO CALLOUT — the hook
% ═══════════════════════════════════
\block{}{
\coloredbox[bgcolor=HeroBGcolor,
            framecolor=HeroBordercolor,
            linewidth=3pt]{
\begin{center}
\large
\textbf{Checker-only UPI--TRM solves 9$\times$9 Sudoku where PPO/A2C/DQN score 0\%.}\\[2pt]
\textbf{33$\times$ lower policy drift at 8$\times$ depth mismatch under inner-loop contraction.}
\end{center}
}
}

% ═══════════════════════════════════
% ARCHITECTURE DIAGRAM — simplified to avoid overlaps
% ═══════════════════════════════════
\block{}{
\vspace{-4mm}
\begin{center}
\begin{tikzpicture}[
    x=1cm, y=1cm,
    mainbox/.style={draw, rounded corners=5pt,
        minimum height=1.2cm,
        font=\small, align=center,
        thick, line width=0.6pt},
    statebox/.style={mainbox, fill=TRMBlue!12,
        draw=TRMBlue!75!black, text width=3.2cm},
    evalbox/.style={mainbox, fill=AccentGreen!12,
        draw=AccentGreen!70!black, text width=6.6cm},
    policybox/.style={mainbox, fill=AccentOrange!12,
        draw=AccentOrange!70!black, text width=5.2cm},
    advbox/.style={mainbox, fill=TRMBlue!8,
        draw=TRMBlue!75!black, text width=4.6cm},
    chip/.style={draw=black!30, fill=gray!6,
        rounded corners=3pt, font=\small, inner sep=4pt},
    arrow/.style={-{Stealth[scale=0.85]}, semithick,
        shorten >=2pt, shorten <=2pt},
    dataarrow/.style={arrow, TRMBlue!80!black},
    feedarrow/.style={arrow, AccentGreen!70!black},
    gradarrow/.style={arrow, AccentOrange!80!black},
]
\node[statebox] (state) at (0,0)
  {\textbf{State / Plan}\\$s_t=(x,y_t)$};
\node[evalbox] (eval) at (8.6,0)
  {\textbf{Checker + Evaluator}\\
   $z^{(t+1)}\!=\!(\Pi_R\!\circ\!f_\theta)(z^{(t)},y,x)$\\
   $U_n\!=\!V_\psi(z^{(n)},x)$;\enspace $r_t\!=\!c(x,y_t)$};
\node[advbox] (adv) at (18.6,0)
  {\textbf{Centered Advantage}\\
   $\hat A\!=\!\hat Q-\E_\pi[\hat Q]$};
\node[policybox] (policy) at (28.4,0)
  {\textbf{Conservative Update}\\
   $\pi_{\mathrm{new}}\!=\!(1\!-\!\alpha)\pi
   \!+\!\alpha\pi_{\mathrm{cand}}$};
\draw[dataarrow] (state.east) --
  node[midway, above=4pt, fill=white,
    inner sep=1pt, font=\small]{$(x,y_t)$}
  (eval.west);
\draw[feedarrow] (eval.east) --
  node[midway, above=4pt, fill=white,
    inner sep=1pt, font=\small]{$r_t,\;U_n$}
  (adv.west);
\draw[gradarrow] (adv.east) --
  node[midway, above=4pt, fill=white,
    inner sep=1pt, font=\small]{$\hat A$}
  (policy.west);
\node[chip, fill=TRMBlue!8] at (8.6,-2.0)
  {\textcolor{TRMBlue}{\textbf{TRM dials:}}
   $L_z,\ n,\ R$};
\node[chip] at (28.4,-2.0)
  {\textbf{Knobs:} $K,\ \alpha$};
\end{tikzpicture}
\end{center}
\vspace{-5mm}
{\small Checker feedback and the $n$-step unrolled evaluator produce the advantage signal; the conservative update edits the plan and repeats.}
}

% ═══════════════════════════════════
% TWO-COLUMN BODY
% ═══════════════════════════════════
\begin{columns}

% ── LEFT ──
\column{0.5}

\block{Motivation}{
\textbf{Tiny Recursive Models (TRMs)} solve reasoning
tasks by iteratively editing a candidate plan---but
require ground-truth supervision.

\medskip
\textbf{Can we train from checker feedback alone?}
(e.g., ``does this Sudoku satisfy constraints?'')

\medskip
\textbf{Challenge:} The evaluator is \textbf{approximate}
(limited capacity) and \textbf{truncated} (finite
depth~$n$). How does truncation interact with policy
improvement?

\medskip
\textbf{Key idea:} Formalize plan editing as a
\textbf{discounted MDP}; analyze via policy iteration
with a compute-truncated value oracle.
}

\block{Plan-Space MDP \& Dials}{
\renewcommand{\arraystretch}{1.15}
\begin{tabular}{@{}ll@{}}
\textbf{State} & $s\!=\!(x,y)$: instance + plan \\
\textbf{Action} & $y'\!=\!\mathrm{edit}(y,a;x)$ \\
\textbf{Reward} & Checker $c(x,y)$ + shaping \\
\end{tabular}

\medskip
\begin{center}
\renewcommand{\arraystretch}{1.2}
\setlength{\tabcolsep}{4pt}
\begin{tabular}{@{}lll@{}}
\toprule
\textbf{Type} & \textbf{Param} & \textbf{Controls} \\
\midrule
\textcolor{TRMBlue}{\textbf{TRM}} & $L_z\!<\!1$
  & Contraction \\
\textcolor{TRMBlue}{\textbf{Dials}} & $n$
  & Unroll depth \\
  & $R$ & Projection radius \\
\midrule
\textcolor{KnobGray}{\textbf{Std}} & $K$
  & Bootstrap horizon \\
\textcolor{KnobGray}{\textbf{Knobs}} & $\alpha$
  & Mixture weight \\
\bottomrule
\end{tabular}
\end{center}
}

\block{Value-Error Decomposition}{
\coloredbox[bgcolor=EqBoxBGcolor,
            framecolor=TRMBlueBorder,
            linewidth=2pt]{
{\small
\textbf{Proposition (simplified).} Under contraction
($L_z\!<\!1$), over the advantage-evaluation closure:
\[
\|U_n \!-\! V^\pi\|_{\infty,\pi}
\le
\underbrace{\frac{\varepsilon_{\mathrm{res}}^*}
  {1\!-\!\gamma^K}}_{\text{architectural}}
+\;
\underbrace{L_V\frac{L_z^n}
  {1\!-\!L_z}C_z}_{\text{truncation}}
\]
}
\smallskip
{\small Truncation bias decays \textbf{geometrically}
as $L_z^n$.}
}
}

\block{Conservative Improvement}{
\coloredbox[bgcolor=EqBoxBGcolor,
            framecolor=OrangeBorder,
            linewidth=2pt]{
{\small
\textbf{Theorem (ideal exact-centering case).} For
$\pi_{\mathrm{new}}\!=\!(1\!-\!\alpha)\pi
\!+\!\alpha\pi_{\mathrm{cand}}$:
\[
\eta(\pi_{\mathrm{new}})
\ge
\widehat{L}_\pi(\pi_{\mathrm{new}})
\!-\!\frac{\alpha\varepsilon_{A,\mathrm{cand}}}
  {1\!-\!\gamma}
\!-\!\frac{2\varepsilon_{\mathrm{CPI}}
  \gamma\alpha^2}{(1\!-\!\gamma)^2}
\]
}
\smallskip
{\small Exact centering makes evaluation error scale with
$\alpha$; with GAE, an added centering-defect term
$\varepsilon_{\mathrm{cent}}$ appears.}
}
}

\makeatletter
\global\advance\TP@blocktop by -12mm
\makeatother

\block{Main Contributions}{
\small
\textbf{1. Checker-only UPI--TRM}\\[-1pt]
Learns plan edits from checker feedback and rewards,
without expert trajectories.\\[6pt]
\textbf{2. Value-error decomposition}\\[-1pt]
Separates architectural residual from geometric
truncation bias $O(L_z^n)$.\\[6pt]
\textbf{3. Conservative improvement}\\[-1pt]
With exact statewise centering, CPI evaluation error is
linear in $\alpha$.
\vspace{8mm}
}

% ── RIGHT ──
\column{0.5}

\block{UPI--TRM Algorithm}{
\coloredbox[bgcolor=GreenLight,
            framecolor=GreenBorder,
            linewidth=1.5pt]{
\textbf{1.\;Value Regression:}
$K$-step bootstrap,
min $(U_n(s)\!-\!\mathrm{sg}(G^{(K)}))^2$
}
\vspace{2mm}
\coloredbox[bgcolor=BlueLight,
            framecolor=TRMBlueBorder,
            linewidth=1.5pt]{
\textbf{2.\;Advantage Estimation:}
$\widehat{A}\!=\!\widehat{Q}
\!-\!\E_\pi[\widehat{Q}]$\enspace
(statewise centered)
}
\vspace{2mm}
\coloredbox[bgcolor=OrangeLight,
            framecolor=OrangeBorder,
            linewidth=1.5pt]{
\textbf{3.\;Conservative Update:}
$\pi_{\mathrm{new}}\!=\!(1\!-\!\alpha)\pi
\!+\!\alpha\pi_{\mathrm{cand}}$,
distill to $\pi_\phi$
}
}

\block{4$\times$4 Sudoku Feasibility}{
\begin{center}
\renewcommand{\arraystretch}{1.3}
\setlength{\tabcolsep}{4pt}
\begin{tabular}{@{}lcc@{}}
\toprule
\textbf{Setting} & \textbf{UPI--TRM}
  & \textbf{PPO/A2C/DQN} \\
\midrule
Easy (1--4 empties) & 90.7--93.3\% & 52\% (random) \\
Hard (6--8 empties) & 48--57\% & \textbf{0\%} \\
\bottomrule
\end{tabular}
\end{center}
\vspace{1mm}
{\small 5k--20k steps, 3 seeds. Ranges denote
contraction-on to no-contraction results; on hard
instances, 11 tuned baseline configs all score 0\%.}
}

\block{Depth-Mismatch Stability}{
$n_{\mathrm{train}}\!=\!2$, eval at $8\times$ mismatch
(both conditions use $R\!=\!10$; projection alone gives
$6$--$10\times$ $\Delta_V$ reduction):

\begin{center}
\renewcommand{\arraystretch}{1.2}
\setlength{\tabcolsep}{4pt}
\begin{tabular}{@{}lccc@{}}
\toprule
& $\Delta_V\!\downarrow$ & $\Delta_\pi\!\downarrow$
& Argmax$\!\uparrow$ \\
\midrule
No Contr. & 0.156 & 0.0063 & 96.0\% \\
Contr.\ ($L_z\!=\!0.9$)
  & \textbf{0.038} & \textbf{0.0002}
  & \textbf{99.0\%} \\
\midrule
\textbf{Gain}
& \textcolor{TRMBlue}{\textbf{4.1$\times$}}
& \textcolor{TRMBlue}{\textbf{33$\times$}}
& --- \\
\bottomrule
\end{tabular}
\end{center}
\vspace{2mm}
\coloredbox[bgcolor=BlueLight,
            framecolor=TRMBlueBorder,
            linewidth=1.2pt]{
\small \textbf{Takeaway:} Contraction lowers drift at
every tested mismatch depth; at $8\times$, it cuts
$\Delta_V$ by \textbf{4.1$\times$} and $\Delta_\pi$ by
\textbf{33$\times$}.
}
}

\block{9$\times$9 Sudoku Results}{
\begin{center}
\renewcommand{\arraystretch}{1.2}
\setlength{\tabcolsep}{4pt}
\begin{tabular}{@{}lccc@{}}
\toprule
\textbf{Method} & \textbf{Success}
& \textbf{Score} & \textbf{$\Delta$} \\
\midrule
\textbf{UPI--TRM}
& \textcolor{TRMBlue}{\textbf{4.0$\pm$4.0\%}}
& \textcolor{TRMBlue}{\textbf{53.3$\pm$1.1}}
& \textcolor{TRMBlue}{\textbf{+26.5}} \\
A2C & 0.0$\pm$0.0\% & 31.9$\pm$0.4 & +5.1 \\
DQN & 0.0$\pm$0.0\% & 29.5$\pm$0.3 & +2.7 \\
PPO & 0.0$\pm$0.0\% & 28.2$\pm$1.5 & +2.2 \\
\bottomrule
\end{tabular}
\end{center}
{\small 81 cells, 729 actions, 50k steps,
mean$\pm$std over 3 seeds.
UPI--TRM is the \textbf{only method solving any puzzles};
baselines cannot propagate constraints across the 81-step horizon.}
}

\block{Takeaways}{
{\small
\textbf{(1)} $L_z$, $n$, $R$ are stability dials:
truncation bias decays as $L_z^n$.
\textbf{(2)} Conservative $\alpha$ limits sensitivity
to imperfect evaluation.
\textbf{(3)} UPI--TRM solves from checker feedback alone;
PPO/A2C/DQN get 0\%.}

\medskip
{\small \textbf{Contact:} buiksat@meta.com}
}

\end{columns}

\end{document}
